\section{Proof of Theorem~\ref{thm:main}}\label{sec:proof}

OpenAI's energy method gives the decay rate $D_0=n_0^2(U-\frac12\log2)=\DecayOA\ldots$ for our parameters.
Section~\ref{sec:signed} raises it to $D=\Decay\ldots$ With $\Lambda=\LamDen\ldots$ the polynomial
$Q_N=(d_0N)!\,\cD_N\Delta_N$ has integer coefficients, degree at most $53N$, and
\[
 |Q_N(x)|\le e^{-\eta N^2}\qquad(|x-G|\le e^{-\rho N}),\qquad \eta=\EtaTot .
\]
Two further properties are needed to choose the prime in the criterion: nonvanishing of $Q_p(a/b)$ when
$p^2\nmid b$ (Proposition~\ref{prop:nonvanishing}), and divisibility of the $i$-th coefficient of $Q_N$ by
$p'^{2i}$ for every prime $N<p'\le2N$ (Lemma~\ref{lem:squares}). We now check the hypotheses of
Proposition~\ref{prop:criterion}.

Take $d=53$, and let the polynomials $Q_p$ of the proposition be the $Q_N$ with $N=p$. Let
\begin{gather*}
 \rho=\Rho,\quad \eta=\EtaTot,\quad \omega=\tfrac35,\quad \delta=2.9\cdot10^{-7},\\
 \lambda=\frac{410}{\rho}=\LamCrit\ldots
\end{gather*}
\begin{enumerate}[label=(\alph*)]
\item $\deg Q_N\le53N$ by \eqref{eq:QN} and Lemma~\ref{lem:degree}.
\item By Proposition~\ref{prop:nonvanishing}, $\Delta_p(a/b)\ne0$ for large $p$ with $p^2\nmid b$. Hence
$Q_p(a/b)\ne0$.
\item By Lemma~\ref{lem:squares}, $p'^{2i}$ divides the $i$-th coefficient of $Q_p$ for all primes
$p<p'\le2p$.
\item By Theorem~\ref{thm:real}, Proposition~\ref{prop:denominators} and $\log(53N)!=o(N^2)$, for
$|x-G|\le e^{-\rho N}$ and large $N$,
\[
 |Q_N(x)|\le\exp\bigl(-(D-\Lambda+o(1))N^2\bigr)\le e^{-\eta N^2},
\]
since $\eta=\EtaTot<D-\Lambda=\EtaExact\ldots$ In the formal proof this is the inequality $\eta<D-\Lambda$ with
$\int\kappa$ replaced by its lower bound $\ChiLean$.
\end{enumerate}
The numerical conditions hold as well. We have $\rho\lambda=410$ and $0<\eta<2d=106$. Next $\omega=\frac35\le1$ and
$2\omega=1.2<\lambda$. Further $2(\omega-\delta(2\lambda-\omega))=1.199998\ldots>1$. Finally
$d+4d\delta\lambda=\PayL\ldots<\eta\lambda=\PayR\ldots$ So Proposition~\ref{prop:criterion} gives
$|G-a/b|>b^{-\rho\lambda}=b^{-410}$ for all integers $a$ and all sufficiently large $b$. This proves
Theorem~\ref{thm:main}.

\section{Remarks on the method}\label{sec:remarks}

\subsection{The three ingredients}
For fixed parameters the method gives the exponent $d\rho/\eta$, and its three ingredients come from
different places. The degree $d=n-g$ comes from the construction. Of $\rho=\Rho$, the prescribed columns
contribute $n_0c_2\approx208.6$ through the energy, and the weight and the binomial factor of one column
contribute $W_0\approx88.0$. Bounding these two factors separately, by their maxima, would add about $5$ to
$\rho$. The number $\eta$ is a small difference of two large numbers, $D\approx6.4\cdot10^3$ and
$\Lambda\approx6.4\cdot10^3$. The corank bounds of Section~\ref{sec:content} contribute about $22$ to it, and the
signed bound of Section~\ref{sec:signed} about $8$. Without them $\eta$ would be about $8$. Since the exponent
is inversely proportional to $\eta$, a further gain of $0.1$ in $\eta$ would lower it by about $1$.

\subsection{The row determinant}
\begin{remark}\label{rem:L}
The energy method bounds $|\det\mathsf L|$ by $6^n$, which is $e^{O(N)}$ and invisible at the scale $N^2$.
At the extremal configuration, numerical evaluation for small $N$ suggests that $\log|\det\mathsf L|$ is about
$-28N^2$. So most of its true size is not captured by our bound either. The bound \eqref{eq:eleven} loses only
$e^{O(N)}$, and Theorem~\ref{thm:heine} is an identity. The loss is in the log-gas bounds of
\S\ref{ssec:loggas}, which use fixed trial measures, and in the requirement that the combined functional be
bounded by its maximum over all configurations. A sharper treatment of the denominator $Z_s(\Pi_+)$ would
raise $\Gamma_0$. These numerical values are not part of the formal proof.
\end{remark}

\subsection{Other parameters}
All the results of Section~\ref{sec:content} are lower bounds for the corank, and the proofs give the kernel
vectors explicitly. The closed formulas, the counting inequality and the families derived from it are
proved for all parameters. So are the other arguments of the paper, except the explicit subspaces of
Appendix~\ref{app:windows} and the trial measures, which belong to $(58,9,3,5,3)$. A different choice of
parameters changes $d$, $\rho$ and $\eta$ together, and we make no claim that ours is optimal. Except for the
pair (F4), the function $\kappa$ uses no relation between different families. So at a prime where intervals
overlap the exact corank may be larger than each bound. An upper bound for the corank is not needed. More
independence between families, or larger explicit subspaces, would increase $\eta$ without changing
anything else.

\section{Formalization}\label{sec:formal}

All results of this paper are formalized in Lean~4 \cite{Lean4} with Mathlib \cite{Mathlib}, on top of OpenAI's
formalization of the irrationality of $G$ \cite{OpenAI2026}. The development is available at
\begin{center}
\FormalizationLink.
\end{center}
Theorem~\ref{thm:main} is the theorem
\begin{center}
{\small\texttt{CatalanIrrationalityMeasure.catalan\_irrationalityMeasure\_410}}
\end{center}
in \lean{Solution.lean}. Its statement, in \lean{Challenge.lean}, uses only Mathlib. Catalan's constant is
written as the series $\sum_j(-1)^j/(2j+1)^2$, and the conclusion is $|G-a/b|>b^{-410}$ for all integers $a$
and all sufficiently large $b$. The modules contain no \lean{sorry} and no \lean{native\_decide}. The command
\lean{\#print axioms} shows that only the three standard axioms \lean{propext}, \lean{Classical.choice} and
\lean{Quot.sound} are used. The repository also contains a configuration \lean{comparator.json} for
\lean{lake comparator} \cite{Comparator}, which checks that the solution file proves the challenge statement
with these axioms only. The prime number theorem is taken from the PNT+ project \cite{PNTPlus}. The parts of
OpenAI's formalization that are used were converted to Lean's module system and ported to Lean
4.35.0-rc2. The repository's \texttt{README} describes the changes and the layout of the modules.

The proof combines conventional formal proofs with numerical certificates. These are the inverses of the base
matrices modulo $1009$, the trial sequences and the covers of Appendices~\ref{app:field} and~\ref{app:trials},
and the subspaces, kernel vectors and interval conditions of Appendix~\ref{app:windows}. The certificates are
exact rational and integer computations checked by the Lean kernel without compiled code. The interval
arithmetic used to generate them \cite{Arb} is not trusted.

\section*{Disclosure of AI use}
The mathematics, the Lean formalization, the certificate generators and the text of this paper were produced
by AI systems at the request of the author. The author posed the problems, directed the work and is
responsible for the paper. The systems were Claude Opus~5.5, Claude Sonnet~5.5 and Claude Haiku~5.5
(Anthropic), running in Claude Code, and GPT-6.1 Sol (OpenAI), running in Codex. The underlying
irrationality proof for Catalan's constant and its Lean formalization are due to OpenAI \cite{OpenAI2026}.

\subsection*{Acknowledgements}
We thank the authors of Mathlib and of the PNT+ project, and OpenAI for making its formalization available.
